\RequirePackage[l2tabu, orthodox]{nag}

\documentclass[%
    pdftex,
    oneside,			% Einseitiger Druck.
    12pt,				% Schriftgroesse
    parskip=half,		% Halbe Zeile Abstand zwischen Absätzen.
    %topmargin = 10pt,	% Abstand Seitenrand (Std:1in) zu Kopfzeile [laut log: unused]
    headheight = 33pt,	% Höhe der Kopfzeile
    %headsep = 30pt,	% Abstand zwischen Kopfzeile und Text Body  [laut log: unused]
    headsepline,		% Linie nach Kopfzeile.
    footsepline,		% Linie vor Fusszeile.
    %footheight = 16pt,	% Höhe der Fusszeile
    abstracton,		% Abstract Überschriften
    DIV=calc,		% Satzspiegel berechnen
    BCOR=8mm,		% Bindekorrektur links: 8mm
    headinclude=false,	% Kopfzeile nicht in den Satzspiegel einbeziehen
    footinclude=false,	% Fußzeile nicht in den Satzspiegel einbeziehen
    listof=totoc,		% Abbildungs-/ Tabellenverzeichnis im Inhaltsverzeichnis darstellen
    toc=bibliography,	% Literaturverzeichnis im Inhaltsverzeichnis darstellen
]{scrreprt}	% Koma-Script report-Klasse, fuer laengere Bachelorarbeiten alternativ auch: scrbook

%%%%%%% Package Includes %%%%%%%
\usepackage{xstring}
\usepackage{ifpdf}
\usepackage{ifluatex}
\usepackage{lastpage}
\usepackage{fancyhdr} % Causing Warnings
\usepackage{pdfpages} % pdf-Seiten einbinden
\usepackage[utf8]{inputenc}
\usepackage[T1]{fontenc}
\usepackage{tikz}
\usepackage{xcolor}
\usepackage[margin=2.5cm,foot=1cm,top=3cm,bottom=3cm]{geometry}	% Seitenränder und Abstände
\usepackage[activate]{microtype} %Zeilenumbruch und mehr
\usepackage{varwidth}
\usepackage[onehalfspacing]{setspace}
\usepackage{makeidx}
\usepackage[autostyle=true,german=quotes]{csquotes}
\usepackage{longtable}
\usepackage{enumitem}	% mehr Optionen bei Aufzählungen
\usepackage{graphicx}
\usepackage{xcolor} 	% für HTML-Notation
\usepackage{float}
\usepackage{array}
\usepackage{calc}		% zum Rechnen (Bildtabelle in Deckblatt)
\usepackage{wrapfig}
\usepackage{pgffor} % für automatische Kapiteldateieinbindung
\usepackage[perpage, hang, multiple, stable]{footmisc} % Fussnoten
\usepackage{acronym}
\usepackage[absolute]{textpos}
\usepackage{scrhack} % in Kombination mit listings-Package kommt es zu Warnings, dieses Paket verhindert die Warnings! Ggf. auskommentieren und die Warnings akzeptieren falls Verzeichnisse nicht so dargestellt werden wie gewünscht
\usepackage{listings} % Code-Listings
\usepackage{color, colortbl}  %Für Highlighten der Tabellenzeilen
\usepackage{amsmath}% http://ctan.org/pkg/amsmath
\usepackage[ngerman]{babel}
\usepackage[font=footnotesize,lofdepth,lotdepth]{subfig}
% eine Kommentarumgebung "k" (Handhabe mit \begin{k}<Kommentartext>\end{k},
% Kommentare werden rot gedruckt). Wird \% vor excludecomment{k} entfernt,
% werden keine Kommentare mehr gedruckt.
\usepackage{comment}

\newcommand{\einstellung}[1]{%
    \expandafter\newcommand\csname #1\endcsname{}
    \expandafter\newcommand\csname setze#1\endcsname[1]{\expandafter\renewcommand\csname#1\endcsname{##1}}
} %Einstellungscommand

\newcommand{\langstr}[1]{\einstellung{lang#1}} % Sprache aus Einstellungen laden

\newcommand{\citem}[1]{\item[\texttt{#1}]} % Code-Item für description-Liste

\newcommand{\todo}[1]{\textit{\textcolor{red}{TODO: #1}}} % Todo-Item

\newcommand{\ladefarben}{%
	\definecolor{LinkColor}{HTML}{00007A}
	\definecolor{ListingBackground}{HTML}{FCF7DE}
} % Farben (Angabe in HTML-Notation mit großen Buchstaben)

%% Programmiersprachen Highlighting (Listings)
\newcommand{\listingsettings}{%
	\lstset{%
		language=C++,			% Standardsprache des Quellcodes
		%numbers=left,			% Zeilennummern links
		%stepnumber=1,			% Jede Zeile nummerieren.
		%numbersep=5pt,			% 5pt Abstand zum Quellcode
		%numberstyle=\tiny,		% Zeichengrösse 'tiny' für die Nummern.
		breaklines=true,		% Zeilen umbrechen wenn notwendig.
		breakautoindent=true,	% Nach dem Zeilenumbruch Zeile einrücken.
		postbreak=\space,		% Bei Leerzeichen umbrechen.
		tabsize=2,				% Tabulatorgrösse 2
		basicstyle=\ttfamily\footnotesize, % Nichtproportionale Schrift, klein für den Quellcode
		showspaces=false,		% Leerzeichen nicht anzeigen.
		showstringspaces=false,	% Leerzeichen auch in Strings ('') nicht anzeigen.
		extendedchars=true,		% Alle Zeichen vom Latin1 Zeichensatz anzeigen.
		captionpos=b,			% sets the caption-position to bottom
		%backgroundcolor=\color{ListingBackground}, % Hintergrundfarbe des Quellcodes setzen.
		xleftmargin=0pt,		% Rand links
		xrightmargin=0pt,		% Rand rechts
		frame=single,			% Rahmen an
		frameround=ffff,
		rulecolor=\color{darkgray},	% Rahmenfarbe
		%fillcolor=\color{ListingBackground},
		keywordstyle=\color[rgb]{0.133,0.133,0.6},
		commentstyle=\color[rgb]{0.133,0.545,0.133},
		stringstyle=\color[rgb]{0.627,0.126,0.941},
    aboveskip=1.5em,
	}
}

%%%%%%%%%%%%%%%%%%%%%%%%%%%%% Kopf-/Fußzeilenwechsel %%%%%%%%%%%%%%%%%%%%%%%%%%%
\setlength{\headheight}{40pt}

\newcommand{\setpagestyle}[1]{%
    \fancypagestyle{plain}{%
    \fancyhf{}
    \fancyhead[L]{\autor}
    \fancyhead[R]{
    \rightmark
    }
    \fancyfoot[L]{
        \noindent{\tiny \langfussz\\
            \begin{tabular*}{16cm}{@{\extracolsep{\fill}}l>{\raggedleft}p{8cm}}
                {\tiny \langstand: \datumAbgabe} & 
                {\tiny \langseite\ \thepage\ \langseitevon\ \pageref*{#1}\vspace{1cm}}\tabularnewline
            \end{tabular*}
        }
    }
    }
    \pagestyle{plain}
}

\newcommand{\setpagestylehead}{%
    \setpagestyle{endOfRomanNumbering}
    \pagenumbering{roman}
}

\newcommand{\setpagestylecontent}{
    \setpagestyle{endOfArabicNumbering}
    \pagenumbering{arabic}
}

\newcommand{\setpagestylefoot}{
    \setpagestyle{LastPage}
    \pagenumbering{Alph}
}

% Flag für die Selbstständigkeitserklärung, Default: true
\newif\ifselbsterkl
\selbsterklfalse

% Flag für das Wasserzeichen auf dem Deckblatt, default: false
\newif\ifwatermark
\watermarkfalse

% Flag für roten Vertraulichkeitspunkt, default: false
\newif\ifreddot
\reddotfalse

% Flag für gelben Vertraulichkeitspunkt, default: false
\newif\ifyellowdot
\yellowdotfalse

% Flag für das Unterschriftenblatt, default: false
\newif\ifunterschriftenblatt
\unterschriftenblattfalse

% Flag für Einfügen der Seitenzahl bei Verweis auf Kapitel/Abschnitt, default: false
\newif\ifrefWithPages
\refWithPagesfalse

% Flag für Einfügen der Abstracts in deutsch und englisch, default: false
\newif\ifbothabstracts
\bothabstractsfalse

% Flag für Einfügen des Abkürzungsverzeichnis
\newif\ifabkverz
\abkverzfalse

% Flag für Einfügen des Abbildungsverzeichnisses
\newif\ifabbverz
\abbverzfalse

% Flag für Einfügen des Formelverzeichnisses
\newif\ifformelverz
\formelverzfalse

% Flag für Einfügen des Formelgroessenverzeichnisses
\newif\ifformelgroeverz
\formelgroeverzfalse 

% Flag für Einfügen des Listingsverzeichnisses
\newif\iflistverz
\listverzfalse

% Flag für Einfügen des Tabellenverzeichnisses
\newif\iftableverz
\tableverzfalse

% Flag für Einfügen des Sperrvermerks
\newif\ifsperrvermerk
\sperrvermerkfalse

% Flag für Einfügen des Abstracts
\newif\ifabstract
\abstractfalse

% Flag für Anhang
\newif\ifappendix
\appendixfalse

% Flag für Literaturverzeichnis
\newif\ifliteratur
\literaturfalse

%Flag für untereiltest Literaturverzeichnis
\newif\ifliteraturteilen
\literaturteilenfalse

% Flag für Glossar
\newif\ifglossar
\glossarfalse

% Flag für Inhaltsverzeichnis
\newif\ifinhalt
\inhaltfalse

% Flag für Reviewer
\newif\ifreviewer
\reviewerfalse


\einstellung{martrikelnr}
\einstellung{titel}
\einstellung{kurs}
\einstellung{datumAnfang}
\einstellung{datumAbgabe}
\einstellung{firma}
\einstellung{firmenort}
\einstellung{abgabeort}
\einstellung{abschluss}
\einstellung{studiengang}
\einstellung{dhbw}
\einstellung{betreuer}
\einstellung{gutachter}
\einstellung{zeitraum}
\einstellung{arbeit}
\einstellung{autor}
\einstellung{sprache}
\einstellung{schriftart}
\einstellung{kapitelabstand}
\einstellung{spaltenabstand}
\einstellung{zeilenabstand}
\einstellung{zitierstil}
\einstellung{selbsterkl}
\einstellung{semester}
\einstellung{studienrichtung}
\einstellung{jahrgang}
\einstellung{abteilung}
\einstellung{standort}
\einstellung{formelbeschreibung} % verfügbare Einstellungen
\documentclass[
    pdftex,
    oneside,			% Einseitiger Druck.
    12pt,				% Schriftgroesse
    parskip=half,		% Halbe Zeile Abstand zwischen Absätzen.
    %topmargin = 10pt,	% Abstand Seitenrand (Std:1in) zu Kopfzeile [laut log: unused]
    headheight = 33pt,	% Höhe der Kopfzeile
    %headsep = 30pt,	% Abstand zwischen Kopfzeile und Text Body  [laut log: unused]
    headsepline,		% Linie nach Kopfzeile.
    footsepline,		% Linie vor Fusszeile.
    %footheight = 16pt,	% Höhe der Fusszeile
    abstracton,		% Abstract Überschriften
    DIV=calc,		% Satzspiegel berechnen
    BCOR=8mm,		% Bindekorrektur links: 8mm
    headinclude=false,	% Kopfzeile nicht in den Satzspiegel einbeziehen
    footinclude=false,	% Fußzeile nicht in den Satzspiegel einbeziehen
    listof=totoc,		% Abbildungs-/ Tabellenverzeichnis im Inhaltsverzeichnis darstellen
    toc=bibliography,	% Literaturverzeichnis im Inhaltsverzeichnis darstellen
]{scrarticle}	% Koma-Script report-Klasse, fuer laengere Bachelorarbeiten alternativ auch: scrbook
 % lese Einstellungen

\newcommand{\iflang}[2]{%
  \IfStrEq{\sprache}{#1}{#2}{}
}

\langstr{abkverz}
\langstr{anhang}
\langstr{glossar}
\langstr{deckblattabschlusshinleitung}
\langstr{artikelstudiengang}
\langstr{studiengang}
\langstr{anderdh}
\langstr{von}
\langstr{dbbearbeitungszeit}
\langstr{dbmatriknr}
\langstr{dbkurs}
\langstr{dbfirma}
\langstr{dbbetreuer}
\langstr{dbgutachter}
\langstr{sperrvermerk}
\langstr{erklaerung}
\langstr{abstract}
\langstr{listingname}
\langstr{listlistingname}
\langstr{listingautorefname}
\langstr{selbsterkl}
\langstr{formelsammlung}
\langstr{kopfz}
\langstr{fussz}
\langstr{seite}
\langstr{seitevon}
\langstr{stand}
\langstr{formelgroeverz} % verfügbare Strings
\input{lang/\sprache} % Übersetzung einlesen

% Einstellung der Sprache des Paketes Babel und der Verzeichnisüberschriften

\PassOptionsToPackage{english, ngerman}{babel}

\iflang{de}{
    \usepackage{babel}
    \selectlanguage{ngerman}
}
\iflang{en}{
    \usepackage{babel}
    \selectlanguage{english}
}
%%%%%% Configuration %%%%%

%% Anwenden der Einstellungen
\usepackage{\schriftart}
\ladefarben

% Titel, Autor und Datum
\title{\titel}
\author{\autor}
\date{\datum}

% PDF Einstellungen
\usepackage[%
    pdftitle={\titel},
    pdfauthor={\autor},
    pdfsubject={\arbeit},
    pdfcreator={pdflatex, LaTeX with KOMA-Script},
    pdfpagemode=UseOutlines, 		% Beim Oeffnen Inhaltsverzeichnis anzeigen
    pdfdisplaydoctitle=true, 		% Dokumenttitel statt Dateiname anzeigen.
    pdflang={\sprache}, 			% Sprache des Dokuments.
]{hyperref}

% (Farb-)einstellungen für die Links im PDF
\hypersetup{%
    colorlinks=true, 		% Aktivieren von farbigen Links im Dokument
    linkcolor=black, 	    % Farbe festlegen
    citecolor=LinkColor,
    filecolor=LinkColor,
    menucolor=LinkColor,
    urlcolor=LinkColor,
    %linktocpage=true, 		% Nicht der Text sondern die Seitenzahlen in Verzeichnissen klickbar
    linktoc=all,            % Seitenzahlen und Text klickbar
    bookmarksnumbered=true 	% Überschriftsnummerierung im PDF Inhalt anzeigen.
}
% Workaround um Fehler in Hyperref, muss hier stehen bleiben
\usepackage{bookmark} %nur ein latex-Durchlauf für die Aktualisierung von Verzeichnissen nötig

% Schriftart in Captions etwas kleiner
\addtokomafont{caption}{\small}

% Literaturverweise (sowohl deutsch als auch englisch)
\iflang{de}{%
\usepackage[
    backend=bibtex,		% empfohlen. Falls biber Probleme macht: bibtex
    bibwarn=true,
    bibencoding=utf8,	% wenn .bib in utf8, sonst ascii
    sortlocale=de_DE,
    style=\zitierstil,
]{biblatex}
}
\iflang{en}{%
\usepackage[
    backend=bibtex,		% empfohlen. Falls biber Probleme macht: bibtex
    bibwarn=true,
    bibencoding=utf8,	% wenn .bib in utf8, sonst ascii
    sortlocale=en_US,
    style=\zitierstil,
]{biblatex}
}
\ladeliteratur{}

% Glossar
\usepackage[nonumberlist,toc]{glossaries}
\usepackage{blindtext} % Blindtext-Package. Common Usage: \blindtext für einzelnen Abschnitt, \Blindtext für mehrere Abschnitte

%%%%%% Additional settings %%%%%%

% Hurenkinder und Schusterjungen verhindern
% http://projekte.dante.de/DanteFAQ/Silbentrennung
\clubpenalty = 10000 % schließt Schusterjungen aus (Seitenumbruch nach der ersten Zeile eines neuen Absatzes)
\widowpenalty = 10000 % schließt Hurenkinder aus (die letzte Zeile eines Absatzes steht auf einer neuen Seite)
\displaywidowpenalty=10000

\setcounter{biburlnumpenalty}{100}
\setcounter{biburlucpenalty}{100}
\setcounter{biburllcpenalty}{100}

% Bildpfad
\graphicspath{{images/}}

% Einige häufig verwendete Sprachen
\lstloadlanguages{PHP, Python, Java, C, C++, bash}
\listingsettings{}
% Umbennung des Listings
\renewcommand\lstlistingname{\langlistingname}
\renewcommand\lstlistlistingname{\langlistlistingname}
\def\lstlistingautorefname{\langlistingautorefname}

% Abstände in Tabellen
\setlength{\tabcolsep}{\spaltenabstand}
\renewcommand{\arraystretch}{\zeilenabstand}

\usepackage{xspace}
\newcommand{\lastcontentpage}{}
\usepackage{amsfonts}

\usetikzlibrary{shapes,arrows,calc,positioning}
% !TeX root = ../dokumentation.tex

%% Some custom TikZ Styles for common UML Elements
%% This could be used instead of TikZ-UML

\tikzstyle{transition} = [%
    thin,
    -angle 45,
    rounded corners
]

\tikzstyle{rect} = [%
    draw,
    rounded corners,
]

\tikzstyle{singlestate} = [%
    draw,
    rounded corners,
    execute at begin node={\begin{varwidth}{10em}}, % Maximum Width
    execute at end node={\end{varwidth}},
    minimum width=10em                              % Minimum Width
]

\tikzstyle{synchronize} = [%
    draw,
    fill=black,
    minimum width=3cm
]

\tikzstyle{decision} = [%
    draw,
    diamond,
    minimum height=0.5cm,
    aspect=1
]

\tikzstyle{tnode} = [%
    font=\footnotesize
]

\tikzstyle{start} = [%
    circle,
    draw,
    fill=black
]

\tikzstyle{final} = [%
    circle,
    double,
    fill=black,
    double distance=0.1cm,
    draw
]

\usepackage{relsize}

\usepackage{censor}

\usepackage{eso-pic}

\usepackage{adjustbox}

%% Paket um Textteile drehen zu können
% \usepackage{rotating}
%% Paket um Seite im Querformat anzuzeigen
%\usepackage{lscape}

\newcommand\Watermark{%
    \put(0,0){%
        \parbox[b][\paperheight]{\paperwidth}{%
            \vfill
            \includepdf[scale=0.8,angle=50,pages={1},pagecommand={}]{ads/watermark}
            \vfill
        }
    }
}

\ifrefWithPages
    %RJG8FE: add a pageref to autoref whenever the referenced page is not the same as the current one
    %        useful for printed documents without clickable hyperlinks
    \AtBeginDocument{\let\oldautoref\autoref}
    \AtBeginDocument{
        \renewcommand{\autoref}[1]{%
            \oldautoref{#1}%
            \ifthenelse{\thepage=\pageref{#1}}% if current page number equals the referenced page number
            {}% then add nothing
            { (S. \pageref{#1})}% else add the text
        }
    }
\fi

\usepackage{amssymb} % Erweiterung der Symbole in Mathematikumgebung

\iflang{de}{\usepackage{icomma}} % Europäsiches Komma in Formeln

\usepackage{tabularx}
\usepackage{tabulary}

%%%%%% Custom Enviroments %%%%%%
\newenvironment{conditions}[1][\formelbeschreibung]
  {%
      #1\tabularx{\textwidth-\widthof{#1}}[t]{
      >{$}l<{$}  @{} >{${}}c<{{}$} @{} X@{}
      }%
  }
  {\endtabularx\\[\belowdisplayskip]}
%%%%%% Custom Commands %%%%%%
\newrobustcmd*{\citecompleteauthor}{\AtNextCite{\DeclareNameAlias{labelname}{given-family}}\citeauthor}